%!TEX root = ../thesis.tex

\chapter{Introduction}\label{chap:introduction}

% Short framing chapter (~5 pages). The substance lives in \cref{chap:background}
% (fuzzing, RL) and \cref{chap:related-work} (RL for fuzzing); this chapter only
% states what the thesis asks, what it contributes, and how it is organized.

\section{Motivation}\label{sec:introduction-motivation}

Modern-day automatic security testing is essential for finding vulnerabilities
and keeping software secure.
Coverage-guided greybox fuzzing plays a major role and has an impressive track
record: OSS-Fuzz alone has helped fix over \num{13000} security vulnerabilities
and \num{50000} bugs across \num{1000} projects~\cite{arya_oss-fuzz_nodate}.
Its success rests on exploring a program by iteratively refining a set of inputs
with the goal of maximizing execution coverage, exposing bugs along the way.
Strikingly, it never forms an explicit model of what makes an input good,
but relies only on a scalar coverage novelty signal, ensuring high execution
throughput.

While this sparse feedback signal enables testing of a large number of different inputs,
it also limits the amount of information that can be extracted from a program's execution trace.
\textcite{bohme_fuzzing_2020} show that while a linearly increasing compute budget
linearly reduces the time to find a given set of bugs, discovering
new bugs requires an exponentially increasing one.

The decision of where to spend the execution budget is largely determined by
hand-tuned heuristics.
This is suboptimal since it has been shown that an optimal fuzzing schedule
is strongly target-dependent~\cite{woo_scheduling_2013,bohme_coverage-based_2016}.
A large part of the research surveyed in \cref{sec:fuzzing-limitations} can be
read as an attempt to extract more information from the same executions.

Stateful targets present an even larger problem.
Coverage-guided fuzzing models a program as a pure function from an
input to a coverage vector, and resets it to its initial state before each new input.
Stateful targets are different: their behavior depends not only on a new input,
but on a (potentially very large) history of inputs.
Examples for such targets include highly relevant software systems such as
industrial control systems, and embedded and IoT
stacks~\cite{zhang_iot_2014,sadeghi_security_2015,borcherding_bandits_2023,ang_quic-fuzz_2025}.

The information available to a fuzzer changes depending on how targets are executed.
While native execution provides high throughput, using emulation allows for easier
program instrumentation without the need for recompilation and enables a fuzzer to
run programs compiled for different architectures.
Some emulators also allow the fuzzer to take a lightweight snapshot of a programs'
execution state and restore it afterwards, making emulation an attractive direction
for stateful fuzzing.

Fuzzing can be modeled as sequential decision making under uncertainty with a scalar progress signal.
Reinforcement learning is a machine learning paradigm that aims at solving these
kinds of problems.
While model-free RL methods have been applied to fuzzing before~\cite{bottinger_deep_2018,binosi_rainfuzz_2023,wang_reinforcement_2021,wang_syzvegas_2021,borcherding_bandits_2023},
the use of model-based RL promises a higher sample efficiency, requiring fewer
executions of a program to achieve a given goal.

EfficientZero~\cite{ye_mastering_2021} is a model-based RL method that learns a
latent model of the environment's dynamics and plans inside it with Monte-Carlo
tree search, reaching above-average human performance on Atari from roughly two
hours of game experience~\cite{kaiser_model-based_2024}.
It allows an agent to learn how to choose an optimal set of actions given
potentially incomplete observations from the environment.
This makes it an interesting research direction for fuzzing.

\section{Problem Statement and Research Questions}\label{sec:introduction-rqs}

The main question this thesis addresses is whether model-based RL can be
successfully used as the core of a fuzzer. EfficientZero is extremely slow
compared to traditional fuzzers, but can learn a useful model of its
environment and make better use of information.

We pose three research questions:

\begin{enumerate}[label=\textbf{RQ\arabic*.}, ref=RQ\arabic*, leftmargin=*]
  \item\label{rq:learnability}
    Can EfficientZero learn a useful policy on fuzzing-shaped rewards?
    Coverage rewards are sparse and observations are aliased and only a partial
    view of a program's true state.

  \item\label{rq:action-space}
    How should the fuzzing action space be modelled?
    We pursue two approaches: modeling actual inputs a program expects as actions,
    as well as modeling mutation operators acting on a corpus as actions.

  \item\label{rq:planning}
    Does planning with a learned model help? A core part of EfficientZero is
    \enquote{imagining} how a program behaves given an input. How well can this be learned?
\end{enumerate}

\section{Contributions}\label{sec:introduction-contributions}

This thesis makes three contributions.

\begin{enumerate}
  \item \emph{Two MDP formulations of emulator-based greybox fuzzing} that allow
    RL algorithms to drive fuzzing of various target executables.

  \item \emph{An implementation}: a reusable EfficientZero implementation in Rust
    using the \code{Burn} deep-learning framework, and the fuzzer \emph{mugiwara}
    built on top of the MOGI emulation framework.

  \item \emph{An empirical evaluation} including a detailed ablation study
    on a simple json parser as well as fuzzing runs for various programs.
\end{enumerate}

\section{Outline}\label{sec:introduction-outline}

The rest of this document is structured as follows.

\cref{chap:background} introduces coverage-guided greybox fuzzing and
reinforcement learning, \cref{chap:efficientzero} describes EfficientZero, and
\cref{chap:related-work} surveys prior applications of RL to fuzzing and the
difficulties they report.

\cref{chap:fuzzing-mdp} models emulator-based fuzzing as a Markov decision
process: it states the interface a model-based learner demands, the constraints
that interface imposes on any fuzzing MDP, and the two different ways we model
the action space.
\cref{chap:implementation} describes the resulting system, from the learning
agent down to the MOGI emulation platform and the coverage and snapshot
primitives it provides.

\cref{chap:targets} introduces the fuzzing targets, and \cref{chap:methodology}
the experimental methodology.
\cref{chap:results} reports the results for both formulations,
\cref{chap:failure-modes} the observed failure modes and \cref{chap:discussion}
compares \emph{mugiwara} against conventional fuzzers.
\cref{chap:conclusion} concludes the research questions and
\cref{chap:future-work} outlines follow-up work.
